\section{Linear temporal completeness and decision: proofs}
\label{sec:annex-p08-linear-temporal-completeness}

\begin{annexitem}{annex:p08.ltl.strict-bridge}{thm:p08-ltl-strict-bridge}{Every endpoint in the strict/non-strict translation}
\begin{proof}
A strict-until witness \(j>n\) is equivalently a witness
\(j\geq n+1\), and its intermediate points \(n<k<j\) are
exactly \(n+1\leq k<j\). This proves the first identity by the
next clause. A non-strict witness is either \(j=n\), giving
\(B\) now, or \(j>n\); the second case additionally requires
\(A\) at \(n\), and its remaining interval is the strict one.
This proves the second identity in both directions. Atoms and
Boolean connectives preserve translation pointwise, and next
preserves it by shifting one position. Structural induction
therefore proves the nested translations. Taking the left until
operand true gives the eventuality identity; Boolean duality gives
invariance. Taking that operand false forces a strict witness to
be the immediate successor on \(\mathbb N\), whereas the
non-strict witness must be current. These last assertions need
the specified discrete successor order.
\end{proof}
\end{annexitem}

\begin{annexitem}{annex:p08.ltl.soundness}{thm:p08-ltl-soundness}{All schemas, both theorem-only rules and suffix invariance}
\begin{proof}
At a fixed word position, uniform substitution into a classical
tautology remains true under the truth values of its substituted
formulas. X-distribution and X-negation follow by evaluating the
indicated formulas at the unique successor. An until witness is
either current, giving \(B\), or later, giving \(A\) now and
the same until at the successor. Conversely either disjunct gives
the required witness. This proves the unfolding equivalence.
MP preserves truth at every position. If a formula is valid at
all positions, it is true at every successor, proving XNec.

For UInd, suppose its displayed premise
\(B\lor(A\land\mathsf XC)\to C\) is valid. If
\(A\mathbin{\mathsf U}B\) holds at \(n\), choose a witness
\(j\geq n\). At \(j\), \(B\) and the valid premise give
\(C\). Working backwards, if \(n\leq k<j\), the until
witness gives \(A\) at \(k\), and the already established
\(C\) at \(k+1\) gives \(\mathsf XC\) there. The valid
premise therefore gives \(C\) at \(k\). Finite backward
induction reaches \(n\). This proves the rule's validity
preservation; an open premise was never assumed globally.
Induction on a finite Hilbert derivation proves soundness.

Finally, for any \(n\), set \(V_n(i)=V(n+i)\). Induction on
the formula proves that truth at \(i\) in this suffix word equals
truth at \(n+i\) in the original. In the until case, subtract
\(n\) from a witness and its intermediate indices, or add it
back for the converse. Thus validity at zero for all words implies
validity at every position; the reverse implication is immediate.
\end{proof}
\end{annexitem}

\begin{annexitem}{annex:p08.ltl.noncompactness}{thm:p08-ltl-noncompactness}{Every finite subset and the failure of finitary strong completeness}
\begin{proof}
If every \(\mathsf X^n p\) holds at a position, \(p\) holds
at each present or later position, exactly the condition for
\(\mathsf Gp\). For a finite subset choose \(N\) at least
every exponent appearing in it, taking \(N=0\) if none appear.
Let \(p\) hold at positions zero through \(N\) and fail at
\(N+1\). All selected premises hold at zero, but \(\mathsf Gp\)
does not. The same word satisfies the finite subset together with
\(\neg\mathsf Gp\). The whole set cannot be satisfied.
Every finitary derivation from premises uses finitely many of them.
If it were both sound and strongly complete, semantic consequence
of \(\mathsf Gp\) from the entire set would produce such a
finite proof, contradicting the finite countermodel. The explicitly
defined finite-premise proof relation has this finite-support
property by definition.
\end{proof}
\end{annexitem}

\begin{annexitem}{annex:p08.ltl.next-boolean}{thm:p08-ltl-next-boolean}{Derived Boolean preservation using the registered rules}
\begin{proof}
Necessitate the theorem \(C\to D\), then apply X-distribution
and MP to derive \(\mathsf XC\to\mathsf XD\).
Necessitating truth proves \(\mathsf X\topv\). X-negation,
with classical double negation, then proves that \(\mathsf X\bottom\)
is false. Negation preservation is the registered schema itself.
Apply implication preservation to \(C\land D\to C\) and
\(C\land D\to D\) for one direction of conjunction preservation.
For the other, necessitate the tautology
\(C\to(D\to C\land D)\), and apply X-distribution twice;
classical propositional reasoning yields
\(\mathsf XC\land\mathsf XD\to\mathsf X(C\land D)\).
Induct on the number of conjuncts, including the empty truth case.
De Morgan's definition of disjunction and negation preservation
give finite disjunction preservation, including the empty falsity
case. Both directions of any proved equivalence follow by applying
implication preservation separately.
\end{proof}
\end{annexitem}

\begin{annexitem}{annex:p08.ltl.type-partition}{thm:p08-ltl-type-partition}{Finite characteristic formulas and consistency}
\begin{proof}
Treat the finitely many formulas in \(C\) as propositional
substitution values. The disjunction of their \(2^m\) complete
signed conjunctions is a substitution instance of a tautology.
For each inconsistent record the calculus proves \(\neg\chi_t\)
by definition. Finitely many uses of propositional reasoning remove
all such disjuncts, proving the consistent-record disjunction.
Different records disagree on some \(D\), so their conjunction
contains both \(D\) and \(\neg D\) and is refutable.

A record assigning one to \(\bottom\), or violating a negation
or conjunction truth condition, has a characteristic formula
refutable by a propositional tautology. Thus every consistent record
satisfies those conditions. If all consistent records assigned
zero to \(\varphi\), their characteristic formulas would each
imply \(\neg\varphi\). The proved partition would then prove
\(\neg\varphi\), contrary to the hypothesis. This proves the
last assertion without a maximal infinite theory or a choice
principle for arbitrary sets of premises.
\end{proof}
\end{annexitem}

\begin{annexitem}{annex:p08.ltl.successor}{thm:p08-ltl-successor}{Seriality and the exact X/U edge clauses}
\begin{proof}
Necessitate the finite partition and use Boolean preservation to
obtain \(\vdash_L\bigvee_{s\in Q_C}\mathsf X\chi_s\).
If a consistent \(t\) had no successor, each
\(\chi_t\land\mathsf X\chi_s\) would be refutable. Combining
these finitely many refutations with the displayed theorem would
refute \(\chi_t\), a contradiction.

On an edge, \(\chi_s\) implies either \(D\) or \(\neg D\)
according to \(s(D)\). Next preserves that implication and
negation. If \(t(\mathsf XD)\ne s(D)\), its characteristic
formula contradicts \(\mathsf X\chi_s\), so no edge exists.
This proves the X clause in both directions.

Put \(U=A\mathbin{\mathsf U}B\in C\). If \(t(B)=1\),
the unfolding axiom forces \(t(U)=1\). If \(t(U)=1\) and
\(t(B)=0\), unfolding makes \(\chi_t\) imply both \(A\)
and \(\mathsf XU\); consistency gives \(t(A)=1\), and the
same next-implication argument forces \(s(U)=1\).
Conversely, if \(t(A)=s(U)=1\) while \(t(U)=0\), unfolding
makes \(\chi_t\) imply \(\neg\mathsf XU\), whereas
\(\mathsf X\chi_s\) implies \(\mathsf XU\). This refutes
the edge. Together these cases are exactly the displayed until
equivalence. They use \(U,A,B\in C\); the formula
\(\mathsf XU\) can occur in the derivation without being a
coordinate of the record.
\end{proof}
\end{annexitem}

\begin{annexitem}{annex:p08.ltl.eventuality-reach}{thm:p08-ltl-eventuality-reach}{The finite reachability formula and UInd}
\begin{proof}
Let \(D\subseteq Q_C\) contain exactly those states from which
there is a finite path to a state with \(B\) true and with \(A\)
true at every earlier state. Define
\(C_D=\bigvee_{s\in D}\chi_s\).
Every record with \(B\) true belongs to \(D\) by the zero-length
path. The finite partition and its signed literals therefore prove
\(B\to C_D\).

If \(t\notin D\) and \(t(A)=1\), there is no edge from
\(t\) to a member of \(D\): such an edge would extend a
witness path and put \(t\) in \(D\). Thus, for each \(s\in D\),
the calculus proves \(\neg(\chi_t\land\mathsf X\chi_s)\).
Finite disjunction preservation of next yields
\(\vdash_L\chi_t\to\neg\mathsf XC_D\).
If instead \(t(A)=0\), its characteristic formula already
implies \(\neg A\). For \(t\in D\), it implies \(C_D\).
Use the partition once more to derive
\(\vdash_L A\land\mathsf XC_D\to C_D\), and combine this
with \(B\to C_D\). The resulting theorem is exactly the
premise required by the registered UInd rule with invariant
\(C_D\). Hence \(\vdash_L U\to C_D\).

If a consistent record with \(U\) true lay outside \(D\),
its characteristic formula would imply \(U\) and, by pairwise
incompatibility, \(\neg C_D\). The last theorem would refute
that characteristic formula. Therefore the record belongs to
\(D\), supplying the stated finite path. Every disjunction and
every collection of refutations used here is finite.
\end{proof}
\end{annexitem}

\begin{annexitem}{annex:p08.ltl.fair-path}{thm:p08-ltl-fair-path}{Scheduling every obligation with genuine progress}
\begin{proof}
List the finitely many until formulas \(U_1,\ldots,U_r\).
Starting from the prescribed record, perform rounds. In a round,
visit these obligations in order. If \(U_j\) is true at the
current endpoint, append a finite witness path supplied by the
reachability lemma; if it is false, append nothing for that slot.
After all slots, append one successor edge, which exists by
seriality. Fix an ordering of the finite graph and always choose
a shortest lexicographically first witness and first available
successor. Ordinary recursion on \(\mathbb N\) now defines
the path, and the extra edge ensures infinitely many positions,
even when every witness has length zero or there are no untils.

Suppose \(U_j=A\mathbin{\mathsf U}B\) is true at a position
\(i\). Until \(B\) occurs, the edge clause forces \(U_j\)
to remain true at the next position. Choose a later scheduled
slot for \(U_j\), which occurs after finitely many finite path
segments. If \(B\) has already occurred, the obligation is
fulfilled. Otherwise \(U_j\) is still true and that slot
appends a finite path ending with \(B\). Thus every obligation
arising at any position is eventually fulfilled.
\end{proof}
\end{annexitem}

\begin{annexitem}{annex:p08.ltl.fair-truth}{thm:p08-ltl-fair-truth}{All formula cases on a realizing record sequence}
\begin{proof}
Induct on the construction of \(D\), proving the assertion at
all positions simultaneously. Atoms hold by the valuation;
falsity, negation and conjunction follow from the Boolean record
conditions and the induction hypothesis. For next, use
\(t_i(\mathsf XE)=t_{i+1}(E)\) and the hypothesis for \(E\).

Let \(D=A\mathbin{\mathsf U}B\). If \(t_i(D)=1\),
eventuality realization gives a least \(j\geq i\) with
\(t_j(B)=1\). At every \(i\leq k<j\), the edge clause
forces \(t_k(A)=1\) and \(t_{k+1}(D)=1\). The hypotheses
for \(A,B\) turn this interval into a semantic until witness.
Conversely, suppose the semantic until has a witness \(j\).
The hypotheses give \(t_j(B)=1\), hence \(t_j(D)=1\) by
its outgoing edge clause; an outgoing edge exists because the
sequence is infinite. For \(k=j-1\) down to \(i\), the
semantic witness gives \(t_k(A)=1\), and the already established
\(t_{k+1}(D)=1\) forces \(t_k(D)=1\). This finite backward
argument proves the reverse implication and rules out every
falsely labelled negative until. No induction on \(D\) itself
at a later position was assumed.
\end{proof}
\end{annexitem}

\begin{annexitem}{annex:p08.ltl.completeness}{thm:p08-linear-temporal-completeness}{A counterword for every underivable formula}
\begin{proof}
Soundness supplies one implication. If \(\nvdash_L A\), then
\(\nvdash_L\neg\neg A\), because classical reasoning would
otherwise yield \(A\). Apply the finite-type construction to
\(\varphi=\neg A\). The partition lemma provides a consistent
record assigning one to \(\neg A\); the scheduling lemma provides
an infinite realizing path from it. Its truth lemma gives a word
falsifying \(A\) at position zero. Thus validity implies provability.

For finite \(\Gamma\), its local semantic consequence is
equivalent, by the Boolean clauses at each position, to validity
of \((\bigwedge\Gamma)\to A\). Apply the theorem just proved.
This establishes the finite-premise claim, including empty
\(\Gamma\), without contradicting the separate arbitrary-premise
noncompactness example.
\end{proof}
\end{annexitem}

\begin{annexitem}{annex:p08.ltl.batch-realization}{thm:p08-ltl-batch-realization}{Both pending-batch directions, including empty obligations}
\begin{proof}
Suppose first that an allowed obligation sequence empties infinitely
often, and \(U=A\mathbin{\mathsf U}B\) is true at position
\(i\). If \(B\) occurs at or after \(i\), the obligation
is realized. Otherwise the local until equation propagates
\(U\) through every later position. Choose \(j\geq i\)
with \(O_j=\varnothing\). Since \(B\) is false,
\(U\in D_{t_j}\subseteq O_{j+1}\). At every subsequent step
it must remain in the batch until \(B\) occurs. A later empty
batch is therefore impossible, a contradiction.

Conversely, suppose every true until is eventually realized.
At an empty batch choose the next batch exactly \(D_t\);
at a nonempty batch choose exactly \(O\setminus H_t\).
These are allowed transitions and introduce no unnecessary
obligations. Each member of a collected batch was a true until
at its collection point. If its right operand does not occur
immediately, the local equation propagates it, and realization
supplies a finite future occurrence. The finitely many members
have a finite maximum witness time, after which that batch is
empty. Thus every empty-batch occurrence is followed by another
after finitely many steps, producing infinitely many. If there
are no untils, the only obligation set is empty at every step,
and both conditions are vacuous.
\end{proof}
\end{annexitem}

\begin{annexitem}{annex:p08.ltl.automaton-language}{thm:p08-ltl-automaton-language}{Acceptance, truth and the source projection}
\begin{proof}
An accepting automaton run supplies locally coherent types
matching the consumed word's atom valuations. Its empty batches
recur, so the batch lemma gives eventuality realization. The
full truth lemma, which required only local coherence and
realization, identifies every subformula bit with its truth on
that word. The initial bit for \(\varphi\) is one, proving
this language inclusion.

Conversely, for a word satisfying \(\varphi\), set every bit
to its actual subformula truth at each position. Boolean
coherence and the X/until edge clauses follow from the semantic
definitions, including the current-witness case. Every true
until has its semantic finite witness. The batch lemma supplies
an accepting obligation run, so these states give an accepting
automaton run on the original word. This proves the other
language inclusion.

In the source's signed closure, a complete Boolean type contains
exactly one of a formula and its single negation. Restriction
to the coordinates \(C\) gives our truth record. Boolean
coherence ensures the inverse assignment, including formulas
already beginning with a negation and elimination of double
negations. The source local transition consumes precisely the
current full type. Projecting that type to its true atoms is
therefore exactly the letter agreement in the displayed
transition; no first-letter shift or empty transition is added.

There are at most \(2^m\) Boolean records and \(2^r\)
obligation sets. All their defining conditions are finite
Boolean tests, so enumeration is effective. Since \(m,r\leq n\),
the stated bound follows. The source's bound using at most
\(2n\) signed closure formulas is looser and is implied by
this coordinate count.
\end{proof}
\end{annexitem}

\begin{annexitem}{annex:p08.ltl.generalized-realization}{thm:p08-ltl-generalized-realization}{The recurrence condition in both directions}
\begin{proof}
Fix \(U=A\mathbin{\mathsf U}B\). Suppose the sequence
realizes eventualities but visits \(F_U\) only finitely often.
Beyond some position its records therefore all make \(U\)
true and \(B\) false. Realization for \(U\) at that position
would give a later \(B\), a contradiction.

Conversely suppose each \(F_U\) recurs, and \(t_i(U)=1\).
If no \(B\) occurs at or after \(i\), the local until
equation forces \(U\) to remain true at every later position.
No such position can belong to \(F_U\), contradicting its
recurrence. Hence \(B\) does occur. This proves realization
for every true until. The case with no untils has two vacuous
conditions and still requires an infinite run.

An accepting run of \(\mathcal G_\varphi\) is locally
coherent, matches all atom letters and has the required initial
bit. Apply the equivalence just proved and the full truth lemma
to obtain truth of \(\varphi\). In the other direction take
actual subformula truth records on a satisfying word. They
satisfy all local equations and realize every until, so the
equivalence gives generalized acceptance. There are at most
\(2^m\) records and exactly the listed \(r\) acceptance
sets. The already proved language characterization of the
pending-batch automaton gives their equality of languages.
\end{proof}
\end{annexitem}

\begin{annexitem}{annex:p08.ltl.generalized-conversion}{thm:p08-ltl-generalized-conversion}{Every round must actually be completed}
\begin{proof}
For \(h\geq1\), an original run determines a unique counter
sequence starting at one. If every \(F_j\) recurs, whenever
the counter waits for \(F_j\), some finite later position
belongs to that set. The counter then advances, or resets after
\(F_h\). Finitely many such waits complete a round. Repeating
this argument shows infinitely many actual visits to
\(F^*\), proving one language inclusion.

For the converse first suppose \(h>1\). Between successive
accepting visits the counter resets to one and can return to
\(h\) only after encountering, in order, states in
\(F_1,\ldots,F_{h-1}\). Each accepting visit itself
witnesses \(F_h\). Infinitely many accepting visits therefore
give infinitely many witnesses to every \(F_j\) in the
projected original run. When \(h=1\), the counter is always
one and \(F^*=F_1\times\{1\}\); acceptance says exactly
that the original run visits \(F_1\) infinitely often. All
transitions preserve the original letter, including the first
one, so projection proves the other inclusion on the same word.

If a required set is empty, no accepting run can complete
infinitely many rounds: the counter becomes stuck before its
empty requirement, or, for the last set, cannot make an
accepting visit. This agrees with generalized acceptance.
An empty start set gives no run in either automaton. If
\(h=0\), generalized acceptance imposes no conditions, and
declaring every original state accepting accepts exactly all
existing infinite runs. The displayed counts follow from the
state products, and substitution of \(|Q_b|\leq2^m\) and
\(h=r\) proves the LTL bound.
\end{proof}
\end{annexitem}

\begin{annexitem}{annex:p08.ltl.single-start}{thm:p08-ltl-single-start}{Preserving the consumed first letter and infinite acceptance}
\begin{proof}
Given an old accepting run \(q_0q_1\cdots\), replace only its
first state by \(q_*\). The first transition on \(a_0\)
exists because \(q_0\in I\) and \(q_1\in\Delta(q_0,a_0)\).
All subsequent transitions are unchanged. Conversely any run
from \(q_*\) chooses a first transition belonging to some
old \(q_0\in I\); restore that state to obtain an old run
consuming the same letters. Adding or removing one initial
accepting visit cannot change infinite recurrence. The new
state is never revisited. If \(I\) is empty its transition
set is empty and neither construction has an infinite run.
\end{proof}
\end{annexitem}

\begin{annexitem}{annex:p08.ltl.accepting-cycle}{thm:p08-ltl-accepting-cycle}{The positive-return criterion and finite lasso bounds}
\begin{proof}
On an accepting run some accepting state \(f\) occurs infinitely
often, because the state set is finite. Its first occurrence
is reachable from a start, and two distinct occurrences give
a path of positive length back to \(f\). Conversely a
prefix to \(f\) followed by infinite repetitions of a
nonempty return walk visits \(f\) infinitely often. Choose
one enabling letter for each edge in these two finite walks;
their repeated letter sequence is an accepted word.

A shortest path from the start set to \(f\) has no repeated
vertex, so its length is at most \(N-1\). A shortest positive
return to \(f\) has no repeated interior vertex, nor an
interior visit to \(f\); otherwise shorten it. Its length is
therefore at most \(N\). This yields the stated prefix/period
bounds. A return walk lies in \(f\)'s strongly connected
component. Conversely, in a cyclic strongly connected component
containing \(f\), paths from \(f\) to a cycle and back give
a positive return. A singleton without a self-loop has no
such return, proving the precise component criterion.
\end{proof}
\end{annexitem}

\begin{annexitem}{annex:p08.ltl.product-correctness}{thm:p08-ltl-product-correctness}{Counterexample paths and accepting product runs in both directions}
\begin{proof}
An infinite product run projects to a legal path from \(w_0\)
and to an automaton run on that path's exact label sequence.
If accepting, the automaton-language theorem makes
\(\neg\varphi\) true on that sequence. Hence the path
violates the proposed universal specification. Conversely any
violating path yields a word satisfying \(\neg\varphi\);
choose its accepting automaton run and pair the corresponding
positions. Every pair is a product vertex and every step a
product edge. Thus it is an accepting product run. Apply the
finite-graph positive-return theorem for the equivalence and
the counterexample prefix/loop. This construction uses the
automaton for the syntactic negation directly, not a presumed
complementation algorithm for nondeterministic automata.
\end{proof}
\end{annexitem}

\begin{annexitem}{annex:p08.ltl.finite-graph-decision}{thm:p08-ltl-finite-graph-decision}{The complete strongly connected component computation}
\begin{proof}
First traverse edges from every initial vertex, marking exactly
the reachable vertices: a discovered vertex has the recorded
path to it, and induction on path length proves no reachable
vertex is missed. Restrict the graph to them. Perform depth-first
search, recording a vertex only when all its successors have
been explored. Then traverse the reversed edges by depth-first
search in decreasing order of these finishing times, skipping
vertices already assigned to an output component.

To prove correctness, contract actual strongly connected
components. The resulting graph is acyclic: a directed cycle
of components would make their vertices mutually reachable.
If there is an edge from component \(C\) to a different
component \(D\), the greatest first-pass finishing time in
\(C\) exceeds that in \(D\). If \(C\) is discovered
first, its search reaches \(D\), which finishes before the
active search in \(C\). If \(D\) is discovered first, it
cannot reach \(C\), so it finishes before \(C\) is later
discovered. The assertion follows in both cases.

The next unassigned component with greatest finishing time
therefore has no incoming edge from an unassigned different
component. In the reversed graph a search from it cannot
leave for another unassigned component, while strong
connectivity makes it reach all of its own vertices.
Remove this component and repeat the argument. Hence every
output component is exact. Test each for an accepting vertex
and for either more than one vertex or a self-loop. The
positive-return theorem proves the resulting decision.
Each pass visits each vertex and edge a bounded number of
times, so it is linear in the explicit graph size.

The formula automaton has \(2^{O(n)}\) states and at most
their square times \(2^{|P|}\) explicit transitions, still
\(2^{O(n)}\). Its local bit tests are polynomial in \(n\).
The product has at most \(|W|2^{O(n)}\) vertices and
\(|R|2^{O(n)}\) edges; label tests have the same bound.
Construction and the two passes therefore have the stated
time bounds. This algorithm is allowed to store the whole
graph; it does not establish the separate space bound.
\end{proof}
\end{annexitem}

\begin{annexitem}{annex:p08.ltl.polynomial-space}{thm:p08-ltl-polynomial-space}{Bounded reachability with an explicit recursion and counters}
\begin{proof}
An automaton state is encoded by \(b=m+r\) bits. Boolean
validity of a type, initial/accepting membership and the
existence of an edge between two encoded states are decidable
in polynomial space by their finite clauses. For satisfiability,
an edge can use the unique letter given by the current type's
atom bits. For the product add a model-state index and check
its input label and transition. The total encoding length is
\(b'=O(n+\log\lVert K\rVert)\).

For valid encoded vertices define a recursive predicate
\(\operatorname{Reach}(u,v,k)\) for paths of length at most
\(2^k\). At \(k=0\), test equality or an edge. At \(k>0\),
enumerate every valid encoded intermediate vertex \(w\),
and test both
\(\operatorname{Reach}(u,w,k-1)\) and
\(\operatorname{Reach}(w,v,k-1)\). Existence of such a
\(w\) is equivalent to the desired path: split any path
after at most \(2^{k-1}\) edges, padding the other half
by zero length if necessary; conversely concatenate the two
paths. This proves the recursion by induction on \(k\).

The number of represented vertices is at most \(2^{b'}\),
and any reachable pair has a simple path of length smaller
than this number. Thus \(k=b'\) suffices. Each recursion
level stores its two endpoints, current intermediate index,
level and loop-control information in \(O(b')\) bits;
the depth is \(O(b')\). Reuse workspace between the two
recursive calls and between consecutive intermediate vertices.
Total workspace is \(O((b')^2)\), and every loop is finite.

Enumerate a start \(i\), an accepting vertex \(f\), and a
vertex \(v\) with an edge \(f\to v\). Test
\(\operatorname{Reach}(i,f,b')\) and
\(\operatorname{Reach}(v,f,b')\). These are exactly the
positive-return conditions, including a self-loop but excluding
an empty return. The outer enumeration adds only linear
workspace. Apply it to \(\mathcal A_\varphi\) for
satisfiability and to the product with
\(\mathcal A_{\neg\varphi}\) for a violating model path.
Negating a halting deterministic answer uses no additional
asymptotic space, giving validity and universal model checking.
The proof describes a terminating deterministic procedure;
no nondeterministic-space conversion theorem is assumed.
\end{proof}
\end{annexitem}

\begin{annexitem}{annex:p08.ltl.pspace-hardness}{thm:p08-ltl-pspace-hardness}{Every field of the configuration reduction}
\begin{proof}
Use a one-tape deterministic decider with a polynomial bound on
the cells it visits, counting its input storage. This is a
permitted definition of deterministic polynomial space. If a
fixed finite number of tapes is used instead, store their symbols
and head markers as tracks on one tape; finite sweeps collect
the scanned symbols in the finite control and then update the
marked cells and head positions. This simulation increases
space only by a constant factor and need not preserve time.
Choose \(s\) with sufficient blank padding and two boundary
cells outside the visited work region. On the fixed input the
actual run never attempts to cross those boundaries.

Here is the promised local function \(f\). In a well-formed
configuration a cell either has an ordinary tape symbol or
contains the unique pair \((q,a)\) marking the head. If the
center is such a pair and the machine transition is
\((q,a)\mapsto(q',d,\varepsilon)\), where
\(\varepsilon\in\{-1,0,1\}\), its new content is \((q',d)\)
when \(\varepsilon=0\), and the plain symbol \(d\) otherwise.
If the center is a plain symbol, it gains a head and the new
state when the head in its left neighbor moves right or the
head in its right neighbor moves left; its own tape symbol is
unchanged in this case. If neither event occurs its content
stays unchanged. These cases are exclusive for a globally
well-formed configuration. Assign any fixed value on malformed
triples, and keep the boundary cells fixed by the separate
boundary constraints. This defines a finite table computable
from the machine. Accepting and rejecting stationary transitions
leave both the marked symbol and all other cells unchanged.

The clock conditions force each block of \(s\) consecutive
positions, starting at a multiple of \(s\), to have phases
\(0,\ldots,s-1\). Exactly-one cell symbols and the unique-head
condition make each block a well-formed configuration, with
the correct fixed boundaries. The initial literals give exactly
the initial input configuration. At phase \(i-1\), the three
tested positions encode the left, center and right cells of
that configuration. Advancing \(s+1\) positions reaches cell
\(i\) of the next configuration. Thus the displayed local
formula forces its correct updated content. Applied to every
interior cell, it forces the entire next configuration.
Induction on blocks identifies every block with the unique
run of \(M\). An eventual accepting-head symbol therefore
implies that \(M\) accepts the input.

Conversely, if \(M\) accepts, list all its configurations
and repeat its accepting configuration forever. Encode their
cells and phases in successive blocks. Every displayed
condition holds: the update formulas use exactly the table
above, and the finite accepting time witnesses the final
eventuality. This gives a satisfying infinite word. If the
machine rejects, its actual run never has an accepting head,
and the forward argument precludes any satisfying word.

For fixed \(M\), \(\mathcal D_M\) and its transition table have
constant size. The width \(s\) is polynomial in the input.
Clock and initial conditions have polynomial size. The
exactly-one-head condition has \(O(s^2)\) pair exclusions,
each using at most \(O(s)\) next symbols. The local-update
conditions have \(O(s)\) positions times the fixed number
of symbol triples, each of size \(O(s)\). Expansion of
Boolean abbreviations and \(\mathsf F,\mathsf G\) preserves
a polynomial bound. Thus the reduction is polynomial-time
and maps any deterministic polynomial-space language to
LTL satisfiability.

The upper bound has already supplied a deterministic
polynomial-space procedure. Flipping the answer of a halting
deterministic decider shows closure of polynomial space
under complement. For any language in that class, apply the
same reduction to a decider for its complement and negate
the resulting formula. It is valid exactly on the original
language. Hence validity is PSPACE-hard as well, and the
proved upper bounds yield both completeness classifications.
\end{proof}
\end{annexitem}

\begin{annexitem}{annex:p08.ltl.model-checking-complete}{thm:p08-ltl-model-checking-complete}{The fixed-clock block translation and the single initial state}
\begin{proof}
Every path after its initial \(w_*\) cycles through phases
\(0,\ldots,k-1\), with an independently selectable bit at
each position. Its blocks therefore define an arbitrary
infinite sequence of valuations of the \(k\) original atoms.
Conversely, any such valuation sequence determines a legal
path by choosing each prescribed bit in its phase. The
single extra start does not restrict the first bit, since
both phase-zero bit states are its successors.

Induct on a formula to compare truth of \(T(A)\) at a
phase-zero position with truth of \(A\) at the corresponding
decoded position. The atomic clause reads exactly its assigned
bit offset. Falsity and both Boolean constructors follow
immediately. Advancing \(k\) steps reaches the next decoded
valuation, proving the X clause. For until, a translated
witness must satisfy \(c\), so it is precisely a current
or later phase-zero position. At every preceding phase-zero
position, the left operand reduces to \(T(A)\); at all
other phases it is automatically true because \(\neg c\)
holds. The induction hypotheses for the two operands now
give exactly the original until witness interval, in both
directions. This includes a witness in the current block.

The initial \(\mathsf X\) moves from \(w_*\) to the first
encoded valuation. Quantifying over all paths therefore
quantifies over all original valuation words. By suffix
invariance, validity at zero is the same as validity at all
positions, proving the reduction claim. There are \(2k+1\)
states and \(4k+2\) edges. Each atomic or next clause of
the translation introduces at most \(k\) next symbols,
and every other clause adds only a constant number of
connectives to its recursively translated operands. Hence
the formula size is \(O(k\operatorname{len}(A))\), and the construction is
polynomial-time. Validity is PSPACE-hard, so this proves
the lower bound for the single-initial-state universal
model-checking problem. Combine with its upper bound.
\end{proof}
\end{annexitem}
